\documentclass[10pt,a4paper]{amsart}
\usepackage[margin=19mm]{geometry}
\usepackage{amsmath,amssymb,mathtools}
\usepackage[scr=rsfs,cal=euler]{mathalfa}
\usepackage{xfrac}
\usepackage[unicode,hidelinks,bookmarks=false]{hyperref}
\newcommand{\ie}{i.e.\ }
\newcommand{\cf}{cf.\ }
\newcommand{\resp}{respectively\ }
\newcommand{\Subsection}{\subsection}
\providecommand{\nobreakdash}{\nobreak\hbox{-}}
\setlength{\emergencystretch}{2em}
\multlinegap0pt
\usepackage{amssymb}
\usepackage{xcolor}
\usepackage{cancel}
\usepackage[normalem]{ulem}
\usepackage[shortlabels]{enumitem}
\newtheorem{proposition}{Proposition}
\title{On simple compact Lie Skew Braces}
\author{Marco Damele, Andrea Loi}
\date{Source: arXiv:2604.22373v1 (CC BY 4.0)}
\begin{document}
\maketitle

\noindent\emph{Source and licence.} On simple compact Lie Skew Braces. Version arXiv:2604.22373v1 (CC BY 4.0), distributed by arXiv under the Creative Commons Attribution 4.0 International licence, http://creativecommons.org/licenses/by/4.0/, as recorded in the arXiv licence metadata for this exact version and shown on the abstract page https://arxiv.org/abs/2604.22373v1. Full licence notice: \texttt{licenses/arxiv-2604.22373-CC-BY-4.0.txt}.\par\smallskip

\noindent\textit{Editorial scope.} Counted unit: the article's proposition prop:integration-morphisms-PLA (source0.tex lines 581--603). The article's complete original proof of it is delivered with it (source0.tex lines 605--731, one \texttt{proof} environment of 127 lines). No mathematics is written, repaired or re-derived: the statement and the proof are the article's own lines, in the article's own notation, and no retained line is substituted or re-flowed.  No further source-local context is needed: every cross-reference of every retained line resolves inside the two delivered spans.  Nothing was omitted from any retained span, so the package carries no omission record.  No retained line contains a citation, so the wrapper carries no bibliography and no bibliography entry.  No retained line contains a figure environment, an image-inclusion command or drawing code, so no image is delivered and the package declares no asset dependency.  Editorial formatting only, in addition: an AMS \texttt{amsart} preamble that restates the article's own declarations and macros used by the retained lines, namely the environments \texttt{proposition} on separate counters, together with the standard packages those lines need.  The retained environments print in this document as Proposition 1 in order of appearance, whereas the article prints the same items with the numbering of its own full document.  The complete native source of the article as distributed in the arXiv e-print for this exact version is bundled unchanged as \texttt{sources/199086/source0.tex}.
\begin{proposition}\label{prop:integration-morphisms-PLA}
Let
\[
(\mathfrak g_1,[\, ,\,]_{\cdot_1},[\, ,\,]_{\circ_1},\triangleright_1), \qquad
(\mathfrak g_2,[\, ,\,]_{\cdot_2},[\, ,\,]_{\circ_2},\triangleright_2)
\]
be integrable post--Lie algebras, and let
\[
(G_1,\cdot_1,\circ_1), \qquad (G_2,\cdot_2,\circ_2)
\]
be simply connected Lie skew braces integrating them. Let
\[
\varphi:\mathfrak g_1 \longrightarrow \mathfrak g_2
\]
be a morphism of post--Lie algebras. Then there exists a unique morphism of Lie skew braces
\[
F:(G_1,\cdot_1,\circ_1)\longrightarrow (G_2,\cdot_2,\circ_2)
\]
such that
\[
dF_e=\varphi.
\]
\end{proposition}

\begin{proof}
Since $\varphi$ preserves $[\,,\,]_{\cdot_1}$ and $[\,,\,]_{\circ_1}$, it integrates uniquely to Lie group
homomorphisms
\[
F:(G_1,\cdot_1)\to (G_2,\cdot_2),
\qquad
E:(G_1,\circ_1)\to (G_2,\circ_2),
\qquad
dF_e=dE_e=\varphi.
\]
We show that $F=E$. Let
\[
\lambda^{G_1}:(G_1,\circ_1)\to \mathrm{Aut}(G_1,\cdot_1),
\qquad
\lambda^{G_2}:(G_2,\circ_2)\to \mathrm{Aut}(G_2,\cdot_2)
\]
be the lambda-actions, and define
\[
\eta:\mathfrak g_1\to \mathrm{Der}(\mathfrak g_2,[\,,\,]_{\cdot_2}),
\qquad
\eta(\xi)(b):=\varphi(\xi)\triangleright_2 b .
\]
It is immediate that $\eta$ is a Lie algebra homomorphism, hence, since $(G_1,\circ_1)$ is simply
connected, it integrates uniquely to a Lie group homomorphism
\[
\widetilde{\lambda}:(G_1,\circ_1)\to \mathrm{Aut}(G_2,\cdot_2)
\qquad\text{with}\qquad
d\widetilde{\lambda}_e=\eta .
\]

Consider
\[
\Theta:(G_1,\circ_1)\to \mathrm{Aff}(G_2,\cdot_2),
\qquad
\Theta(g):=(F(g),\widetilde{\lambda}_g).
\]
Assume for the moment that $\Theta$ is a Lie group homomorphism. Let
\[
\iota_{G_2}:(G_2,\circ_2)\to \mathrm{Aff}(G_2,\cdot_2),
\qquad
\iota_{G_2}(u):=(u,\lambda^{G_2}_u).
\]
Then $\iota_{G_2}\circ E$ is also a Lie group homomorphism, and
\[
d\Theta_e=d(\iota_{G_2}\circ E)_e,
\]
because the first components are both $\varphi$, while on the second component
\[
d\widetilde{\lambda}_e=\eta=d\lambda^{G_2}_e\circ \varphi=d(\lambda^{G_2}\circ E)_e.
\]
By uniqueness of integration from the simply connected group $(G_1,\circ_1)$,
\[
\Theta=\iota_{G_2}\circ E.
\]
Comparing first components gives $F=E$, so $F$ is a homomorphism for both laws. It remains to prove that $\Theta$ is a homomorphism. For this it suffices to show that
\[
F\circ \lambda^{G_1}_g=\widetilde{\lambda}_g\circ F
\qquad\forall g\in G_1,
\]
for then
\[
\Theta(g)\Theta(h)
=
\bigl(F(g)\cdot_2 \widetilde{\lambda}_g(F(h)),\,\widetilde{\lambda}_g\widetilde{\lambda}_h\bigr)
=
\bigl(F(g\cdot_1 \lambda^{G_1}_g(h)),\,\widetilde{\lambda}_{g\circ_1 h}\bigr)
=
\Theta(g\circ_1 h).
\]

To prove the intertwining relation, it is enough to compare differentials at the identity. Set
\[
\rho_{G_1}:=D_{G_1}\circ \lambda^{G_1}:(G_1,\circ_1)\to \mathrm{Aut}(\mathfrak g_1,[\,,\,]_{\cdot_1}),
\qquad
\widetilde{\rho}:=D_{G_2}\circ \widetilde{\lambda}:(G_1,\circ_1)\to \mathrm{Aut}(\mathfrak g_2,[\,,\,]_{\cdot_2}),
\]
where $D_{G_1}(\alpha)=d\alpha_e$ and $D_{G_2}(\beta)=d\beta_e$. Then
\[
d(\rho_{G_1})_e(\xi)(\zeta)=\xi\triangleright_1 \zeta,
\qquad
d(\widetilde{\rho})_e(\xi)(b)=\varphi(\xi)\triangleright_2 b.
\]
Thus it is enough to prove that
\[
\varphi\circ \rho_{G_1}(g)=\widetilde{\rho}(g)\circ \varphi
\qquad\forall g\in G_1.
\]

Define
\[
\Psi(g)(T):=\widetilde{\rho}(g)\circ T\circ \rho_{G_1}(g)^{-1},
\qquad
T\in \mathrm{Hom}(\mathfrak g_1,\mathfrak g_2).
\]
Then $\Psi$ is a Lie group homomorphism. Hence the map
\[
f(g):=\Psi(g)(\varphi)
\]
is constant as soon as its differential at the identity vanishes. For $\xi,\zeta\in \mathfrak g_1$,
\[
\bigl(d\Psi_e(\xi)(\varphi)\bigr)(\zeta)
=
d(\widetilde{\rho})_e(\xi)(\varphi(\zeta))
-
\varphi\bigl(d(\rho_{G_1})_e(\xi)(\zeta)\bigr)
=
\varphi(\xi)\triangleright_2 \varphi(\zeta)-\varphi(\xi\triangleright_1 \zeta)=0,
\]
since $\varphi$ preserves the post-Lie product. Therefore $f$ is constant, and evaluating at $e$
gives
\[
\varphi\circ \rho_{G_1}(g)=\widetilde{\rho}(g)\circ \varphi
\qquad\forall g\in G_1.
\]

It follows that
\[
d(F\circ \lambda^{G_1}_g)_e=d(\widetilde{\lambda}_g\circ F)_e.
\]
Since $(G_1,\cdot_1)$ is simply connected, uniqueness of integration yields
\[
F\circ \lambda^{G_1}_g=\widetilde{\lambda}_g\circ F
\qquad\forall g\in G_1.
\]
Hence $\Theta$ is a Lie group homomorphism, and the conclusion follows.
Uniqueness is immediate from uniqueness of integration.
\end{proof}

\end{document}
